\begin{table}[H]
\centering
\caption{Niveles de remuneración laboral por departamento, 2025}
\label{tab:dept_remuneracion_niveles_33}
\scriptsize
\begin{tabular}{lrrrrr}
\toprule
Departamento & Rem./trab. 2025 & Puesto & Rem./hora 2025 & Puesto & Brecha \\
\midrule
Bogotá D.C. & 2,89 & 1 & 14,9 & 1 & 0,0\% \\
Putumayo & 2,37 & 2 & 12,4 & 2 & 18,0\% \\
San Andrés y Providencia & 2,26 & 3 & 11,2 & 5 & 22,0\% \\
Vaupés & 2,25 & 4 & 11,8 & 3 & 22,3\% \\
Antioquia & 2,18 & 5 & 11,2 & 4 & 24,6\% \\
Casanare & 2,09 & 6 & 10,3 & 8 & 27,7\% \\
Guainía & 2,06 & 7 & 10,4 & 7 & 28,8\% \\
Cundinamarca & 1,99 & 8 & 10,4 & 6 & 31,4\% \\
Vichada & 1,98 & 9 & 9,9 & 10 & 31,7\% \\
Santander & 1,94 & 10 & 9,8 & 11 & 33,1\% \\
Valle del Cauca & 1,91 & 11 & 10,0 & 9 & 34,1\% \\
Guaviare & 1,87 & 12 & 9,7 & 12 & 35,4\% \\
Risaralda & 1,82 & 13 & 9,4 & 14 & 37,0\% \\
Amazonas & 1,82 & 14 & 9,2 & 16 & 37,2\% \\
Meta & 1,80 & 15 & 8,9 & 18 & 37,9\% \\
Caldas & 1,78 & 16 & 9,4 & 13 & 38,3\% \\
Arauca & 1,77 & 17 & 8,7 & 19 & 38,9\% \\
Quindío & 1,71 & 18 & 9,2 & 15 & 41,0\% \\
Atlántico & 1,61 & 19 & 8,5 & 20 & 44,2\% \\
Tolima & 1,58 & 20 & 8,2 & 21 & 45,5\% \\
Caquetá & 1,56 & 21 & 7,8 & 22 & 45,9\% \\
Boyacá & 1,56 & 22 & 9,0 & 17 & 46,0\% \\
Norte de Santander & 1,50 & 23 & 7,6 & 23 & 48,0\% \\
Huila & 1,40 & 24 & 7,4 & 24 & 51,5\% \\
Cesar & 1,38 & 25 & 7,1 & 26 & 52,3\% \\
Bolívar & 1,36 & 26 & 7,2 & 25 & 52,9\% \\
Magdalena & 1,23 & 27 & 6,5 & 29 & 57,6\% \\
Chocó & 1,22 & 28 & 6,4 & 30 & 57,9\% \\
Cauca & 1,15 & 29 & 6,9 & 27 & 60,2\% \\
Nariño & 1,10 & 30 & 6,8 & 28 & 61,9\% \\
Córdoba & 1,08 & 31 & 6,3 & 31 & 62,6\% \\
La Guajira & 1,02 & 32 & 5,4 & 33 & 64,9\% \\
Sucre & 1,01 & 33 & 6,0 & 32 & 65,2\% \\
\bottomrule
\end{tabular}
\caption*{\footnotesize Nota: remuneración por trabajador en millones de pesos al mes; remuneración por hora en miles de pesos. La brecha se calcula como la distancia porcentual frente al departamento líder en remuneración por trabajador. Un valor de 0\% corresponde al líder; un valor de 30\% indica que el departamento está 30\% por debajo del líder. Fuente: cálculos propios con GEIH.}
\end{table}
